%*******************************************************
% Abstract in German
%*******************************************************
\begin{otherlanguage}{ngerman}
	\pdfbookmark[0]{Zusammenfassung}{Zusammenfassung}
	\chapter*{Zusammenfassung}
Diese Arbeit untersuchte den sprachspezifischen Performance-Overhead von OpenTelemetry in AWS 
Lambda für Java, Python und Node.js. 
Basierend auf fünf Microbenchmarks wurden instrumentierte und nicht-instrumentierte Funktionen in 
einer kontrollierten Umgebung verglichen. 
Die Performance-Metriken Laufzeit, Speicherverbrauch und Initialisierungszeiten wurden gemessen und 
statistisch ausgewertet.

Die Ergebnisse zeigen, dass Python den geringsten Overhead aufweist, 
gefolgt von Java und Node.js. Node.js verursacht insbesondere bei der Init Duration extrem hohe Overheads. 
In der absoluten Performance ist Python in nahezu allen Szenarien die schnellste Sprache. 
Eine Kostenanalyse bestätigt, dass Python in 4 von 5 Benchmarks die günstigste Option ist, 
während Java nur bei der asynchronen Lambdaverkettung wirtschaftlicher ist.

Für die Praxis empfiehlt sich Python für die meisten synchronen Anwendungsfälle, während Java bei 
asynchronen Aufrufen von Lambda konkurrenzfähig ist.

\end{otherlanguage}
